\section{Contractions into \texorpdfstring{$\ell_1$}{ell1}}
\label{sec:contraction}

The small jumps in every chart force small average variation in $\ell_1$.
Throughout this section, scalar gradients carry the Euclidean norm.

\begin{proposition}[Contraction bound]
\label{prop:contraction}
Let $(V,d)$ and the chart maps $v_s:U\to V$ be those of
\Cref{prop:metric}, where $U=(-2,2)^m$, $1\le s\le S=m^3$, and
$N=m^6$. There is an absolute constant $C$ such that every map
$F:V\to\ell_1^D$, for every finite $D$, satisfying
\[
  \|F(v)-F(w)\|_1\le d(v,w)\qquad(v,w\in V)
\]
obeys
\begin{equation}
  \E_{\theta,\eta\in T}
  \|F(v_1(\theta))-F(v_1(\eta))\|_1
  \le C\sqrt m\,(\log m)^{5/2},
  \label{con:eq:contraction}
\end{equation}
where $T=(-1,1)^m$ and the two parameters are independent and uniform.
\end{proposition}

After smoothing, each gradient is a combination of chart normals. The
jump bounds give a total coefficient budget $Cm(\log m)^2$, while the
good-chart property contributes the factor $C\sqrt{(\log m)/m}$.
We then use the following scalar inequality to pass from gradients to
average distances without a further dimension factor.

\begin{lemma}[First-moment inequality on a cube]
\label{con:lem:cube}
Let $T=(-1,1)^m$, and let $h$ be a real-valued $C^1$ function on a
neighborhood of $\overline T$. If $\theta,\eta$ are independent and uniform
on $T$, then
\begin{equation}
  \E|h(\theta)-h(\eta)|
  \le \E|\nabla h(\theta)|.
  \label{con:eq:cube}
\end{equation}
\end{lemma}

We use the Gaussian rotation argument of Maurey and Pisier~\cite{Pis86},
followed by the coordinatewise Gaussian distribution function.

\begin{proof}
Write $\Phi$ for the standard Gaussian distribution function, set
$\psi(t)=2\Phi(t)-1$, and let $\Psi:\R^m\to T$ apply $\psi$
coordinatewise. If $G,G'$ are independent standard Gaussian vectors in
$\R^m$, then $\Psi(G),\Psi(G')$ are independent and uniform on $T$.
For $0\le t\le\pi/2$, put
\[
  Y(t)=\cos(t)G+\sin(t)G',
  \qquad
  Y'(t)=-\sin(t)G+\cos(t)G'.
\]
At each fixed time, $Y(t)$ and $Y'(t)$ are independent standard Gaussian
vectors. Consequently, conditioning on $Y(t)$ gives
\begin{align*}
  \E\left[\left|\frac{d}{dt}h(\Psi(Y(t)))\right|
      \,\middle|\,Y(t)\right]
  &=\sqrt{\frac{2}{\pi}}\,
    \left|D\Psi(Y(t))^{T}\nabla h(\Psi(Y(t)))\right|\\
  &\le \frac{2}{\pi}\,|\nabla h(\Psi(Y(t)))|,
\end{align*}
because $\sup|\psi'|=\sqrt{2/\pi}$. The fundamental theorem of calculus,
followed by Tonelli's theorem, bounds the expected difference between the
endpoints by the integral of the expected absolute derivative. The
distribution of $\Psi(Y(t))$ is uniform on $T$ for every $t$, so integration
over $[0,\pi/2]$ proves \eqref{con:eq:cube}. All the integrands are
integrable: $\nabla h$ is bounded on $\overline T$, and Gaussian vectors have
finite first moments.
\end{proof}

\begin{proof}[Proof of \Cref{prop:contraction}]
We use $r=m^{-8}$ and $\tau=m^{-2000}$ as in \Cref{prop:metric}.
For each chart, define
\[
  f_s(\theta)=F(v_s(\theta))\quad\text{for almost every }\theta\in U,
\]
and extend $f_s$ by zero outside $U$. Set $f_s=0$ also on the threshold hyperplanes and on $\partial U$.
These are bounded, compactly supported functions because $V$ is finite. The
macroscopic estimate in \Cref{prop:metric} implies
\begin{align}
  \|f_s(\theta)-f_1(\theta)\|_1&\le Cr,
       &&\text{for almost every }\theta\in U,\label{con:eq:chart-close}\\
  \|f_s(\theta)-f_s(\eta)\|_1
       &\le c_*\sqrt m\,|\theta-\eta|+Cr,
       &&\text{off the threshold hyperplanes in }U.
       \label{con:eq:macro-variation}
\end{align}

\paragraph{Smoothing.}
Fix an even, nonnegative function $\rho\in C_c^\infty((-1,1))$ with
$\int_\R\rho=1$. Set $\lambda=m^{-2}$ and
\[
  \kappa_\lambda(z)=\prod_{j=1}^m\lambda^{-1}\rho(z_j/\lambda),
  \qquad H_s=f_s*\kappa_\lambda.
\]
The functions $H_s:\R^m\to\R^D$ are smooth. For sufficiently large $m$,
the support of $x\mapsto\kappa_\lambda(\theta-x)$ lies in $U$ whenever
$\theta\in\overline T$. Also
\[
  \int_{\R^m}|\partial_j\kappa_\lambda|
       =\lambda^{-1}\|\rho'\|_{L_1(\R)}.
\]
Every displacement in the support of $\kappa_\lambda$ has Euclidean norm
at most $\sqrt m\lambda$. Thus \eqref{con:eq:macro-variation} gives
\begin{equation}
  \|H_s(\theta)-f_s(\theta)\|_1\le C(m\lambda+r)
  \quad\text{for almost every }\theta\in T.
  \label{con:eq:smoothing-error}
\end{equation}
Since $f_s-f_1$ vanishes outside $U$, \eqref{con:eq:chart-close} bounds its
$\ell_1$ norm by $Cr$ almost everywhere on $\R^m$. Differentiating the
convolution therefore yields, for every $\theta\in T$ and every $j$,
\begin{equation}
  \|\partial_jH_s(\theta)-\partial_jH_1(\theta)\|_1
      \le Cr/\lambda.
  \label{con:eq:derivative-mismatch}
\end{equation}

\paragraph{The derivative carried by each family of interfaces.}
Fix a chart $s$ and abbreviate $u_i=u_{s,i}$. Its interfaces are contained
in the hyperplanes
\[
  Q_{i,k}=\{x\in\R^m:\langle u_i,x\rangle=k\tau\},
  \qquad 1\le i\le N,\quad k\in\mathbb Z.
\]
Only finitely many of these hyperplanes meet $U$. At almost every point
of $Q_{i,k}\cap U$, let $\Delta_{i,k}(x)\in\R^D$ denote the value of
$f_s$ on the side of increasing $\langle u_i,x\rangle$ minus its value
on the opposite side. This is well-defined away from intersections with
the other interface hyperplanes. By the nonparallelism in
\Cref{lem:directions}, those intersections have zero
$(m-1)$-dimensional measure. The two labels at such an interface differ
by $\tau e_i$, so the adjacent-label estimate in \Cref{prop:metric} gives
\begin{equation}
  \|\Delta_{i,k}(x)\|_1
      \le C\tau\frac{m}{N}(\log m)^2
  \quad\text{for almost every }x\in Q_{i,k}\cap U.
  \label{con:eq:jump-bound}
\end{equation}

For completeness, the distributional gradient of each scalar component
of $f_s$ on $U$ is the vector-valued measure
\begin{equation}
  D f_{s,\alpha}
   =\sum_{i,k}\frac{u_i}{|u_i|}\,
       \Delta_{i,k,\alpha}\,
       \mathcal H^{m-1}\!\restriction(Q_{i,k}\cap U).
  \label{con:eq:distributional-gradient}
\end{equation}
Indeed, test against a smooth function compactly supported in $U$ and
integrate by parts on each polyhedral cell. The cell interiors contribute
nothing because $f_s$ is constant there. The two contributions on each
shared facet combine into the jump times the unit normal $u_i/|u_i|$;
lower-dimensional intersections have zero facet measure. The test
function vanishes near $\partial U$, so there is no contribution from
the outer boundary.

For $\theta\in T$, the kernel $\kappa_\lambda(\theta-\cdot)$ is a smooth
test function compactly supported in $U$. Convolving
\eqref{con:eq:distributional-gradient} with this kernel gives, for every
component $\alpha$,
\begin{equation}
  \nabla H_{s,\alpha}(\theta)
       =\sum_{i=1}^N u_{s,i}\,b_{s,i,\alpha}(\theta),
  \label{con:eq:gradient-representation}
\end{equation}
where the vector coefficient $b_{s,i}(\theta)\in\R^D$ is
\begin{equation}
  b_{s,i}(\theta)=\frac{1}{|u_{s,i}|}
       \sum_{k\in\mathbb Z}\int_{Q_{i,k}\cap U}
       \kappa_\lambda(\theta-x)\Delta_{i,k}(x)
       \,d\mathcal H^{m-1}(x).
  \label{con:eq:coefficient-formula}
\end{equation}
This identity holds at every $\theta\in T$: distributional convolution
agrees with the classical derivative of the smooth function $H_s$.

We claim that the coefficients satisfy the pointwise budget
\begin{equation}
  \sum_{\alpha=1}^D|b_{s,i,\alpha}(\theta)|
       \le C\frac{m}{N}(\log m)^2
       \qquad(\theta\in T).
  \label{con:eq:coefficient-budget}
\end{equation}
To establish it, fix $u=u_{s,i}$ and $\theta$, and let $Z$ have density
$x\mapsto\kappa_\lambda(\theta-x)$. Write $h$ for the density of the
real random variable $\langle u,Z\rangle$. The coarea formula for this
linear functional states that
\[
  h(t)=\frac{1}{|u|}
       \int_{\langle u,x\rangle=t}
       \kappa_\lambda(\theta-x)\,d\mathcal H^{m-1}(x).
\]
The lower bound $|u|\ge1/2$ in \Cref{lem:directions} supplies a coordinate
$j$ with $|u_j|\ge1/(2\sqrt m)$. Conditional on the other coordinates of
$Z$, the variable $\langle u,Z\rangle$ has a translated, possibly
reflected copy of the one-dimensional density $\rho$, scaled by
$\lambda|u_j|$. Its derivative has $L_1$ norm
$\|\rho'\|_{L_1}/(\lambda|u_j|)$. Taking the mixture over the other
coordinates shows that $h$ is continuously differentiable, compactly
supported, and satisfies
\begin{equation}
  \int_\R|h'(t)|\,dt\le C\sqrt m/\lambda.
  \label{con:eq:density-variation}
\end{equation}

For any nonnegative, continuously differentiable, compactly supported
density $h$, comparison on each interval $[k\tau,(k+1)\tau]$ gives
\[
  \tau h(k\tau)
     \le\int_{k\tau}^{(k+1)\tau}h(t)\,dt
       +\tau\int_{k\tau}^{(k+1)\tau}|h'(t)|\,dt.
\]
Summing and using \eqref{con:eq:density-variation}, we obtain
\begin{equation}
  \tau\sum_{k\in\mathbb Z}h(k\tau)
       \le1+C\tau\sqrt m/\lambda\le2
  \label{con:eq:lattice-bound}
\end{equation}
for all sufficiently large $m$. Because the convolution kernel is
nonnegative and is supported inside $U$, the triangle inequality in
\eqref{con:eq:coefficient-formula}, followed by
\eqref{con:eq:jump-bound} and \eqref{con:eq:lattice-bound}, proves
\eqref{con:eq:coefficient-budget}.

\paragraph{Combining the charts.}
Fix $\theta\in T$ and a component $\alpha$, and put
$a_\alpha=\nabla H_{1,\alpha}(\theta)$. If $a_\alpha\ne0$, the
direction property in \Cref{lem:directions} gives at least $S/2$ charts
$s$ such that
\[
  \max_i\left|\left\langle
       \frac{a_\alpha}{|a_\alpha|},u_{s,i}\right\rangle\right|
       \le C\sqrt{\frac{\log m}{m}}.
\]
For each such chart, \eqref{con:eq:gradient-representation} implies
\[
  |a_\alpha|
    \le C\sqrt{\frac{\log m}{m}}\sum_i|b_{s,i,\alpha}(\theta)|
       +|\nabla H_{s,\alpha}(\theta)-a_\alpha|.
\]
All terms on the right are nonnegative. Summing over the good charts and
then enlarging that sum to all charts therefore gives
\begin{equation}
  |a_\alpha|\le\frac{2}{S}\sum_{s=1}^S
       \left(
       C\sqrt{\frac{\log m}{m}}\sum_i|b_{s,i,\alpha}(\theta)|
       +|\nabla H_{s,\alpha}(\theta)-a_\alpha|
       \right).
  \label{con:eq:good-chart-average}
\end{equation}
For $a_\alpha=0$ the same inequality holds trivially. The good charts may
depend on $\alpha$ and $\theta$; \eqref{con:eq:good-chart-average} removes
that dependence before we sum over components.

For each $s$, the Euclidean norm is at most the sum of absolute
coordinates, so \eqref{con:eq:derivative-mismatch} gives
\[
  \sum_{\alpha=1}^D
       |\nabla H_{s,\alpha}(\theta)-\nabla H_{1,\alpha}(\theta)|
  \le\sum_{j=1}^m
       \|\partial_jH_s(\theta)-\partial_jH_1(\theta)\|_1
  \le Cmr/\lambda.
\]
Sum \eqref{con:eq:good-chart-average} over $\alpha$ and apply this bound
and \eqref{con:eq:coefficient-budget}. We obtain at every $\theta\in T$
\begin{align}
  \sum_{\alpha=1}^D|\nabla H_{1,\alpha}(\theta)|
   &\le C\sqrt{\frac{\log m}{m}}\,
          N\frac{m}{N}(\log m)^2+Cmr/\lambda\notag\\
   &\le C\sqrt m\,(\log m)^{5/2},
   \label{con:eq:total-gradient}
\end{align}
where $mr/\lambda=m^{-5}$.

Finally, apply \Cref{con:lem:cube} to every component of $H_1$ and sum.
Equation \eqref{con:eq:total-gradient} yields
\[
  \E_{\theta,\eta\in T}\|H_1(\theta)-H_1(\eta)\|_1
       \le C\sqrt m\,(\log m)^{5/2}.
\]
The expected distance for $f_1$ exceeds this by at most
$2C(m\lambda+r)$, by \eqref{con:eq:smoothing-error}. Since
$m\lambda=m^{-1}$, this additional term is absorbed into the same bound.
The definition of $f_1$ now proves \eqref{con:eq:contraction}.
\end{proof}
